\documentclass[11pt,reqno]{amsart}
\usepackage[T1]{fontenc}
\usepackage[utf8]{inputenc}
\usepackage{lmodern,amsmath,amssymb,amsthm,mathtools,booktabs,longtable,array,microtype,xurl,mathrsfs}
\usepackage[a4paper,margin=25mm]{geometry}
\usepackage[hidelinks]{hyperref}
\newtheorem{theorem}{Theorem}[section]
\newtheorem{proposition}[theorem]{Proposition}
\newtheorem{lemma}[theorem]{Lemma}
\newtheorem{corollary}[theorem]{Corollary}
\theoremstyle{definition}\newtheorem{definition}[theorem]{Definition}
\theoremstyle{remark}\newtheorem{remark}[theorem]{Remark}
\newcommand{\Q}{\mathbb Q}\newcommand{\R}{\mathbb R}\newcommand{\C}{\mathbb C}\newcommand{\Z}{\mathbb Z}\newcommand{\OO}{\mathbb O}
\DeclareMathOperator{\tr}{tr}\DeclareMathOperator{\Tr}{Tr}\DeclareMathOperator{\Ann}{Ann}\DeclareMathOperator{\Fix}{Fix}\DeclareMathOperator{\diag}{diag}\DeclareMathOperator{\lcm}{lcm}\DeclareMathOperator{\Ree}{Re}\DeclareMathOperator{\Jac}{Jac}\DeclareMathOperator{\Pic}{Pic}
\newcommand{\eps}{\varepsilon}
\emergencystretch=2em
\title[Three-slot cubes and marked Jordan arithmetic]{Unifying three-slot slices:\\ associator recovery, integral cubes, infinite return,\\ and marked Jordan arithmetic}
\author{Research continuation for The Clankers}
\date{11 September 2026}
\begin{document}
\begin{abstract}
A three-coordinate vector algebra and a genuinely three-slot tensor are different
objects. We retain the distinction and construct explicit maps between the source
associator, three-slot cubes, graded cubic-norm operators, octonionic coordinates,
and marked Erd\H{o}s--Straus arithmetic. Three marked pairwise faces leave an
exact eight-dimensional joint fibre. A concrete integral $3\times3\times3$ cube
has three identical smooth determinant cubics but an infinite-order round trip
through their kernel maps. First-Tits coordinates give three-graded transformations
mixing all blocks and changing the quadratic trace while preserving cubic norm.
All 196560 minimal vectors in Wilson's octonionic Leech construction are replayed,
including the distribution of their retained ternary Jordan determinant. A marked
Erd\H{o}s--Straus solution is embedded reversibly in an integral cube and then in
the infinite-return example, retaining every change of basis and scale. Terzo's
exponential-ring quotient supplies a separate coefficient/period layer; its
unconditional period subgroup is distinguished from the complete kernel identified
under Schanuel's conjecture. No universal Erd\H{o}s--Straus existence theorem,
Schanuel proof, or historical-priority claim is made.
\end{abstract}
\maketitle
\tableofcontents

\section{Scope, conventions, and what is inherited}

The requested object is treated as a three-slot construction, not as a claim that
writing a real vector space as $\R\oplus\C$ has settled its multiplication,
monodromy, or integral arithmetic. Three slots, three spectral values, real
dimension, algebraic degree, and coefficient transcendence degree remain different
marked quantities. No slice is intrinsically designated ``the real part.''

The supplied archive \texttt{New folder.zip} contains an explicit three-input
associator programme and a Terzo exponential-ring programme. The relevant local
claims are read from the supplied mathematical notes and reconstructed below;
private conversations and the original thesis PDF are not redistributed. The
source ledger records exactly which mathematical documents were inspected.

The general genus-one interpretation of cubes is inherited from Bhargava and Ho
\cite[\S\S3.2.2--3.2.3, 5.1.1]{BH}; the first-Tits and exceptional grading
framework is classical, with the explicit coordinate presentation in
\cite[\S\S5--6 and Appendix A]{BP}. Wilson's Leech construction is imported at
its stated scale \cite[\S\S2--4]{Wilson}. The numerical cubes, marked attachments,
and joint certificate calculations here are explicitly derived; they are not
advertised as historically new theorems. Primary-source searching identified
these antecedents, not a comprehensive priority classification.

All matrices act on column vectors. Tensor slots are numbered $0,1,2$; entries
inside a displayed $3\times3$ matrix are numbered $1,2,3$ unless stated otherwise.
Projective rational vectors are represented in certificates by primitive integer
triples whose first nonzero entry is positive. Every inverse below includes its
exceptional locus. A retained inverse does not imply that an arbitrary projection
is injective.

\section{The source's ternary operation, before contraction}

Let $k$ have characteristic different from two, let $V$ have a symmetric bilinear
form $B$, and define the spin product on $k\oplus V$ by
\[
 (a,u)(b,v)=(ab+B(u,v),av+bu).
\]
Use the associator convention $[x,y,z]=(xy)z-x(yz)$.

\begin{proposition}[Exact ternary recovery]\label{prop:recovery}
For $x=(a,u),y=(b,v),z=(c,w)$,
\[
 [x,y,z]=(0,B(u,v)w-B(v,w)u).
\]
If $u\ne0$ and $B(u,u)\ne0$, the operator
\[
 \mathcal R_u(z)=\frac{[(0,u),(0,u),z]}{B(u,u)}
 =\left(0,w-\frac{B(u,w)}{B(u,u)}u\right)
\]
is a projection and
\[
 \Fix(\mathcal R_u)=\operatorname{im}(\mathcal R_u)
 =\Ann(L_{(0,u)})=\{(0,w):B(u,w)=0\}.
\]
The nonisotropy condition is the precise domain of this normalization.
\end{proposition}
\begin{proof}
Expanding both parenthesizations cancels their scalar components and all scalar
multiples of $u,v,w$ except the two displayed terms. Multiplication by $(0,u)$
is $(c,w)\mapsto(B(u,w),cu)$. Its kernel therefore has $c=0$ and $B(u,w)=0$.
The displayed operator kills the scalar part and the $u$ component and is the
identity on that kernel. This proves idempotence and both equalities.
\end{proof}

For the archive's normalized blade $J=\langle1,u,v\rangle$, the form is
$B=-\langle\ ,\ \rangle$, so $u^2=v^2=-1$, $uv=0$, and its complete associator
has four nonzero basis cells:
\[
 [u,u,v]=-v,\quad [u,v,v]=u,\quad [v,u,u]=v,\quad [v,v,u]=-u.
\]
In particular, $-[u,u,v]=v$ and $-[v,v,u]=u$. These statements already occur in
the supplied ternary-recovery source; the continuation is not credited with
originating them. A contraction using only $\alpha(v,u,-)$ and $\alpha(u,v,-)$
is a noninjective linear map on the space of all trilinear tensors. For example,
a tensor supported only on $(u,u,v)$ is invisible to it. Extra knowledge of the
full multiplication can of course recover information that a contraction alone
has discarded.
The source associator is vector-valued: its 27 input cells are vectors, not
27 scalars. Its exact scalar interface is the whole linear family
$\ell\mapsto\ell\circ\alpha$ for $\ell\in J^*$. For a basis $b_r$ with dual
$b_r^*$ the inverse is $\alpha=\sum_r b_r(b_r^*\circ\alpha)$. Retaining the
whole family is lossless; retaining only one scalar component is a further
projection and is not called an isomorphism here.

The original unnormalized seed calculation is also independently replayed. With
Cayley--Dickson recursion
\[
 (a,b)(c,d)=(ac-\bar d b,da+b\bar c),
\]
left-associated blade words, and $C_2$ tags, the archive's seeds in the
32-dimensional tagged sedenion space are
\[
 X=o_1[0]-o_{1234}[1],\qquad Y=o_2[0]+o_{34}[1].
\]
The verifier reconstructs the entire basis multiplication, verifies
$XY=YX=0$, $X^2=Y^2=-2$, computes both annihilator bases by rational row
reduction, and verifies
\[
 \Ann(L_X)=\Fix(-\tfrac12[X,X,-]),\quad
 \Ann(L_Y)=\Fix(-\tfrac12[Y,Y,-]).
\]
Both spaces have dimension eight in each case; both containments, not only one,
are checked on exact bases. This is a finite linear-algebra certificate for the
displayed seeds, not a universal claim about all sedenion zero divisors.

\subsection{The local M\"obius return is an annihilator line}
Take $V=\R^2$ and $B$ negative definite. For
\[
 u_\theta=(\cos(\theta/2),\sin(\theta/2)),\qquad
 v_\theta=(-\sin(\theta/2),\cos(\theta/2)),
\]
Proposition~\ref{prop:recovery} projects onto $\R v_\theta$. This line and its
projection return at $\theta+2\pi$, while $v_{\theta+2\pi}=-v_\theta$.
Trivializing over $[0,2\pi]$ therefore identifies $(2\pi,a)$ with $(0,-a)$:
the annihilator lines form a M\"obius line bundle. A nowhere-zero section would
be a continuous nonzero real coefficient with opposite endpoint signs, which
is impossible by the intermediate value theorem. This is an actual bundle
calculation from the ternary operator. It keeps its local sign return separate
from the all-integer return constructed in Section~\ref{sec:moore}.

\section{Three marked faces do not determine a cube}

Let $V_i$ be three rank-three free modules over a commutative ring, with marked
primitive vectors $e_i$ and functionals $\ell_i$ satisfying $\ell_i(e_i)=1$.
A cube means a trilinear form
\[
 T:V_0\times V_1\times V_2\longrightarrow k.
\]
Its three marked faces are $F_0=T(e_0,-,-)$, $F_1=T(-,e_1,-)$,
$F_2=T(-,-,e_2)$. This is not the collection of all nine coordinate slices.

\begin{theorem}[Complete three-face gluing fibre]\label{thm:glue}
Faces $F_i$ come from a common $T$ if and only if their common edges agree:
\begin{align*}
 g_{01}(w)&=F_0(e_1,w)=F_1(e_0,w),\\
 g_{02}(v)&=F_0(v,e_2)=F_2(e_0,v),\\
 g_{12}(u)&=F_1(u,e_2)=F_2(u,e_1).
\end{align*}
Their common value at all three marks is denoted $c$. A canonical lift is
\begin{align}
 T_0(u,v,w)={}&\ell_0(u)F_0(v,w)+\ell_1(v)F_1(u,w)+\ell_2(w)F_2(u,v)\notag\\
 &-\ell_0(u)\ell_1(v)g_{01}(w)-\ell_0(u)\ell_2(w)g_{02}(v)\notag\\
 &-\ell_1(v)\ell_2(w)g_{12}(u)+c\ell_0(u)\ell_1(v)\ell_2(w).\label{eq:glue}
\end{align}
The complete fibre is
\[
 T_0+W_0\otimes W_1\otimes W_2,\qquad W_i=\{f\in V_i^*:f(e_i)=0\}.
\]
Over a field the face restriction has rank 19 and its kernel has dimension
$2\cdot2\cdot2=8$. Over $\Z$ with these primitive marks the same statement is
an exact split statement of free abelian groups, with no unmentioned rational
lifting condition.
\end{theorem}
\begin{proof}
Necessity follows by evaluating the same trilinear form in two orders.
Substituting $u=e_0$ in \eqref{eq:glue} cancels all terms except $F_0$;
substitution in either other slot does the same. For the kernel, use
$V_i=ke_i\oplus\ker\ell_i$. A trilinear form vanishing when any slot is its
mark is exactly a form on the product of the three complementary modules.
Their dual ranks are two. This also gives the inverse: subtract the canonical
lift and restrict to the three complements. All steps are integral when the
chosen splittings are integral.
\end{proof}

Thus there are eight genuinely joint coefficients invisible to these three
marked pairwise views. This eight is not automatically one of the octonionic
Peirce spaces: equal dimensions do not supply a map between their different
roles. The explicit exceptional-algebra map is given later.

\section{An integral cube with infinite three-slice return}\label{sec:moore}

For the following cube, every slot is a copy of $\Q^3$, with ordered bases kept
separate. Its first pencil is
\[
 M_0(x,y,z)=
 \begin{pmatrix}x&2z&3y\\3z&y&2x\\2y&3x&z\end{pmatrix}.
\]
Equivalently, its three first-axis arrays are
\begin{equation}\label{eq:cube}
 T_0=\begin{pmatrix}1&0&0\\0&0&2\\0&3&0\end{pmatrix},\quad
 T_1=\begin{pmatrix}0&0&3\\0&1&0\\2&0&0\end{pmatrix},\quad
 T_2=\begin{pmatrix}0&2&0\\3&0&0\\0&0&1\end{pmatrix}.
\end{equation}
For any slot $i$, contract that slot and put the remaining two slots, in
increasing order, on rows and columns. Direct determinant expansion gives,
for each of the three pencils,
\begin{equation}\label{eq:hesse}
 \det M_i(x,y,z)=-6F(x,y,z),\qquad
 F=x^3+y^3+z^3-6xyz.
\end{equation}
The verifier checks all three identities as formal polynomials.

\begin{lemma}[Smoothness and exact kernel maps]\label{lem:kernel}
The three determinant curves $C_i$ over $\Q$ are smooth genus-one curves.
At every point of $C_i$, the contracted matrix has rank exactly two. For
$i\ne j$, the incidence condition
\[
 T(v_i,v_j,-)=0
\]
with the slots put in their original order defines an isomorphism
$\phi_{ij}:C_i\to C_j$. Its inverse is $\phi_{ji}$. A local formula is any
nonzero appropriate row or column of the adjugate; a zero chosen cofactor
only requires changing that chart.
\end{lemma}
\begin{proof}
A singular point would satisfy $x^2=2yz$, $y^2=2xz$, $z^2=2xy$.
A zero coordinate forces all coordinates zero. Otherwise multiplication gives
$1=8$, a contradiction in characteristic zero. A smooth plane cubic has genus
one. If a contracted matrix on the curve had rank at most one, all its
cofactors would vanish; differentiation of its determinant would make the
point singular, contrary to what was just proved.

A rank-two matrix has a unique projective left kernel and a unique projective
right kernel. Use the left one for a row slot and the right one for a column
slot. The resulting vector in slot $j$ gives a nonzero kernel vector in the
corresponding pencil $M_j$, so it lies on $C_j$. Its pencil is again rank two,
and the reverse incidence recovers the original point uniquely. Adjugate
charts are regular and agree on overlaps because they span the same unique
kernel line. This proves the global inverse statement. This is the explicit
incidence construction of \cite[\S\S3.2.2, 5.1.1]{BH}.
\end{proof}

Keep the full round trip
\[
 H=\phi_{20}\phi_{12}\phi_{01}:C_0\longrightarrow C_0.
\]
These maps are not to be identified merely because the three curve equations
in \eqref{eq:hesse} look identical.

\begin{theorem}[Faithful all-integer return]\label{thm:infinite}
For the cube \eqref{eq:cube}, $H$ has infinite order. In particular,
$n\mapsto H^n$ is an injective homomorphism $\Z\to\operatorname{Aut}(C_0)$.
\end{theorem}
\begin{proof}
Take $O=[1:-1:0]$ as origin. The triangle-of-isomorphisms theorem
\cite[\S3.2.3, p.~19; \S5.1.1 following Lemma~5.3, p.~35]{BH} identifies $H$
as translation by a rational point $P$ on the elliptic curve $(C_0,O)$.
Its hypotheses hold: the base field is $\Q$ and the cube is nondegenerate by
Lemma~\ref{lem:kernel}. The direction of the three kernel maps is the direction
in that theorem; reversing all maps gives $-P$ and does not change order.

The following exact incidence chain is obtained by row/column cross products:
\[
 (1,-1,0)\xrightarrow{0\to1}(2,3,1)
 \xrightarrow{1\to2}(21,-52,19)
 \xrightarrow{2\to0}(1817,5275,3258).
\]
Nine steps give
\begin{align*}
 H^3(O)=\big(&666161010410184261713443501009747,\\
 &214674741825814609600406026499153,\\
 &518752964649250744517842481308050\big),
\end{align*}
which is not $O$. All intermediate points and inverse tests are in the JSON.

The primitive cubic $F$ has smooth reduction at 5 and 17: the same gradient
argument works since these primes divide neither 3 nor 7. Exhaustive projective
counts, with the first nonzero coordinate set to one, give
\[
 \#C(\mathbb F_5)=9,\qquad \#C(\mathbb F_{17})=21.
\]
The lists of all points are included. Good-reduction injectivity on torsion
prime to the residue characteristic follows from
\cite[Chapter II, Theorem~4.1 and Corollary~4.2, pp.~62--64]{Milne}.
For completeness, the reduction kernel has a separated filtration with
successive groups isomorphic to the additive residue field. An element of
order prime to that characteristic cannot have a nonzero first image in this
filtration. Hence it is zero. Applying this at both 5 and 17 shows that the
order of any rational torsion point divides $\gcd(9,21)=3$: the 5-primary and
17-primary parts are also excluded by reduction at the other prime. If $P$
were torsion, $3P=0$, contradicting $H^3(O)\ne O$. Thus $P$ and $H$ have
infinite order. This argument is not extrapolation from a finite orbit.
\end{proof}

The result is a three-slot tensor with identical plane-cubic
shadows and a retained return invisible to identifying those shadows.
The return is stronger than a parity sign: it distinguishes every integer.

\section{Three-graded operators and the full Albert coordinates}

Take three oriented three-dimensional spaces $V_0,V_1,V_2$. For a typed
connection to a trilinear cube on input spaces $W_0,W_1,W_2$, take
$V_i=W_i^*$ and retain the dual bases and volume forms. Then the trilinear
coefficient tensor is literally in $V_0\otimes V_1\otimes V_2$, not identified
with its dual by an unrecorded inner product. Let
\[
 \mathcal J=\operatorname{Hom}(V_0,V_1)\oplus
 \operatorname{Hom}(V_1,V_2)\oplus\operatorname{Hom}(V_2,V_0).
\]
Write its elements as $(A,B,C)$ with the displayed forward arrows. Bases give
three $3\times3$ blocks, hence 27 scalar entries, and define
\begin{equation}\label{eq:titsnorm}
 N(A,B,C)=\det A+\det B+\det C-\tr(CBA).
\end{equation}
The mixed term traverses the complete directed triangle. It cannot be evaluated
from three unrelated determinant values. The ordering in \eqref{eq:titsnorm}
is forced by column-vector composition.

This is the first-Tits cubic norm, with the standard block order changed by
$(A_0,A_1,A_2)=(A,C,B)$. The basepoint $(I,0,0)$ needed to recover a Jordan
product is an extra marking; it does not assert an intrinsic ``real block.''
The associated trace and adjoint in this order are
\begin{align*}
 \Tr(A,B,C)&=\tr A,\\
 (A,B,C)^\#&=(A^\#-CB,\ C^\#-BA,\ B^\#-AC).
\end{align*}
With $x\#y=(x+y)^\#-x^\#-y^\#$ and
$T(x,y)=\Tr x\Tr y-\Tr(x\#y)$, the product is
\[
 x\circ y=\tfrac12\{x\#y+\Tr(x)y+\Tr(y)x
 -(\Tr(x)\Tr(y)-T(x,y))1\}.
\]
The classical first-Tits theorem supplies its Jordan-algebra status; the
coordinate presentation and exact norm/adjoint conventions are
\cite[\S5, equations (10)--(15)]{BP}.

\subsection{A cube acts on the three blocks}
For a pure tensor $u=a\otimes b\otimes c\in V_0\otimes V_1\otimes V_2$, use
oriented cross products and define
\begin{equation}\label{eq:Du}
 D_u(A,B,C)=
 \big(b(a\times Cc)^T,\ c(b\times Aa)^T,\ a(c\times Bb)^T\big).
\end{equation}
Extend linearly in all 27 tensor entries. Thus a three-slot cube and a
three-block state have different roles: there is an explicit action
$T\mapsto D_T\in\operatorname{End}(\mathcal J)$, rather than an identification
based only on equal dimensions. This action is injective with an explicit
coefficient inverse. Write $T=\sum t_{ijk}e_i\otimes e_j\otimes e_k$ and let
$\epsilon_{nim}$ be the orientation sign. For any distinct $n,i,m$,
\begin{equation}\label{eq:Dinverse}
 t_{ijk}=\epsilon_{nim}
 \big(D_T(0,0,E_{mk})_A\big)_{jn}.
\end{equation}
Indeed the displayed matrix entry of the first component of $D_T$ is
$\sum_l t_{ljk}\epsilon_{nlm}$, and the distinct-index condition leaves only
$l=i$. This proves the inverse over $\Z$ as well as over fields. The checker
verifies all 27 coefficient identities as formal linear polynomials.

\begin{theorem}[Integral simultaneous three-block moves]\label{thm:move}
For pure $u$, one has $D_u^2=0$, and
\[
 U_u(t)=I+tD_u,\qquad U_u(t)^{-1}=I-tD_u,
 \qquad N(U_u(t)X)=N(X).
\]
Over $\Z$, integral $a,b,c,t$ give an integral automorphism of the 27-entry
lattice. In general these maps do not fix the basepoint or preserve the
quadratic trace.
\end{theorem}
\begin{proof}
Applying \eqref{eq:Du} twice produces zero: for example
$\delta C\,c=a(c\times Bb)^Tc=0$, and the other two components are cyclically
the same. This proves the inverse.

The coordinate-change group $\mathrm{SL}(V_0)\times\mathrm{SL}(V_1)\times
\mathrm{SL}(V_2)$ acts by
\[
 A\mapsto g_1Ag_0^{-1},\quad B\mapsto g_2Bg_1^{-1},\quad
 C\mapsto g_0Cg_2^{-1}.
\]
It preserves \eqref{eq:titsnorm}, and the defining volume-form expression for
\eqref{eq:Du} shows equivariance. When all three vectors are nonzero, volume-one
basis changes reduce them to the first basis vector. When one is zero the
operator is zero. It suffices over $\Q$, therefore, to check that basis case;
the resulting integer polynomial identity then holds over every commutative
coefficient ring.

In that case the only changes are
\begin{align*}
 A_{12}&\mapsto A_{12}-tC_{31},& A_{13}&\mapsto A_{13}+tC_{21},\\
 B_{12}&\mapsto B_{12}-tA_{31},& B_{13}&\mapsto B_{13}+tA_{21},\\
 C_{12}&\mapsto C_{12}-tB_{31},& C_{13}&\mapsto C_{13}+tB_{21}.
\end{align*}
Each determinant is affine in $t$, since only its first row changes. Each
increment is $e_1r$ with $r e_1=0$, so products of two increments vanish.
Any quadratic trace-product term also vanishes by cycling two increments
adjacent under trace. The linear derivative satisfies, by a two-minor expansion,
\[
 \delta\det A=\tr(C\,\delta B\,A),\quad
 \delta\det B=\tr(\delta C\,B\,A),\quad
 \delta\det C=\tr(CB\,\delta A).
\]
For the first equality, both sides are
\[
 -C_{31}A_{23}A_{31}+C_{31}A_{21}A_{33}
 +C_{21}A_{21}A_{32}-C_{21}A_{22}A_{31}.
\]
The other equalities follow by cycling $A,B,C$. Thus the entire norm difference
is zero. All coefficients in the inverse and transformation are integral.
\end{proof}

For arbitrary $T$, $D_T^2$ need not be zero. One must not extend the truncated
exponential formula from pure tensors to their sums. Products of pure moves
are the exact finite, invertible transformations used here.

The action lies in the classical $\Z/3$-graded exceptional algebra
\[
 \mathfrak g_0=\mathfrak{sl}(V_0)\oplus\mathfrak{sl}(V_1)\oplus\mathfrak{sl}(V_2),
 \quad \mathfrak g_1=V_0\otimes V_1\otimes V_2,
 \quad \mathfrak g_2=V_0^*\otimes V_1^*\otimes V_2^*,
\]
with dimensions $24,27,27$; the resulting algebra is the classical
$\mathfrak e_6$ representation \cite[\S6, equations (16)--(20)]{BP}.
Independently of that identification, \texttt{graded\_lie.py} constructs the
27 degree-one matrices, extracts 27 independent degree-two commutators,
constructs all 24 trace-free degree-zero matrices, and checks all 3003 unordered
basis-pair commutators in their prescribed grade. Mixed commutators span all
24 degree-zero directions. Jacobi follows from matrix commutators. Dimension
alone is not used as a recognition theorem.

\subsection{A spectrum-changing calculation}
Let
\[
 A=\diag(2,3,5),\quad B=0,\quad C=E_{31},\quad
 (a,b,c)=(e_2,e_1,e_1).
\]
Then the move in Theorem~\ref{thm:move} gives
\[
 A(t)=\diag(2+t,3,5),\quad B(t)=3tE_{13},\quad C(t)=E_{31}.
\]
Define the quadratic trace coefficient
\[
 S(A,B,C)=\tr(A^\#)-\tr(CB).
\]
Direct substitution gives
\[
 N=30,\qquad S=31+5t,\qquad\Tr=10+t.
\]
Thus $4N-pS=120-p(31+5t)$ changes. The earlier obstruction for
unit-fixing Jordan automorphisms does not apply to these larger transformations.
Setting $t=(120/p-31)/5$ makes this intrinsic invariant equation zero, but this
rational split-Albert point has not thereby acquired positive integral spectral
coordinates. The equality is recorded at its actual scope, without treating a
matrix-state solution as an integer-denominator solution.

\subsection{Explicit return to all three octonionic entries}
Use Zorn coordinates $q=(\alpha,u,v,\beta)$ with multiplication
\begin{align*}
 (\alpha,u,v,\beta)(\gamma,w,r,\delta)
 =(&\alpha\gamma+u\cdot r,
     \alpha w+\delta u-v\times r,\\
   &\gamma v+\beta r+u\times w,
     v\cdot w+\beta\delta).
\end{align*}
The norm is $n(q)=\alpha\beta-u\cdot v$, conjugation is
$q^*=(\beta,-u,-v,\alpha)$, and trace is $\alpha+\beta$.

For $X=(A,B,C)$, put $\xi_i=A_{ii}$ and
\begin{align*}
 x_1&=(A_{23},\operatorname{col}_1B,-\operatorname{row}_1C,A_{32}),\\
 x_2&=(A_{31},\operatorname{col}_2B,-\operatorname{row}_2C,A_{13}),\\
 x_3&=(A_{12},\operatorname{col}_3B,-\operatorname{row}_3C,A_{21}).
\end{align*}
Then
\begin{equation}\label{eq:phi}
 \Phi(X)=\begin{pmatrix}\xi_1&x_3&x_2^*\\x_3^*&\xi_2&x_1\\x_2&x_1^*&\xi_3\end{pmatrix}.
\end{equation}
Its inverse reads the diagonal and six scalar corners into $A$, the upper
three-vectors into columns of $B$, and the negatives of the lower three-vectors
into rows of $C$. Hence all 27 entries are retained, with singleton fibres.

The cubic norm of the image is
\[
 \xi_1\xi_2\xi_3-\xi_1n(x_1)-\xi_2n(x_2)-\xi_3n(x_3)
 +\tr((x_1x_2)x_3)=N(X).
\]
One coordinate proof expands the final term as
\[
 A_{12}A_{23}A_{31}+A_{21}A_{32}A_{13}
 -\sum_{i\ne j}A_{ij}(CB)_{ji}+\det C+\det B.
\]
The diagonal-norm terms contribute the other determinant terms and
$-\sum_iA_{ii}(CB)_{ii}$, proving the identity. The complete 45-monomial
identity is also checked symbolically by the standard-library verifier. This
is the explicit first-Tits/Zorn map of \cite[Theorem~11 and Appendix A]{BP},
with its block order changed as stated, not a new claim of discovering it.

Over $\R$ these Zorn coordinates describe split octonions, not the division
real form. Over $\C$, the compact real form is retained by
\[
 \sigma(A,B,C)=(A^\dagger,C^\dagger,B^\dagger).
\]
Indeed its fixed points have $A=A^\dagger$, $C=B^\dagger$, and the trace
quadratic form is $\tr(A^2)+2\tr(BB^\dagger)>0$ off zero. On each Zorn entry,
the corresponding conjugate-linear involution is
\[
 (\alpha,u,v,\beta)\mapsto
 (\bar\beta,-\bar v,-\bar u,\bar\alpha),
\]
whose fixed norm is $|\alpha|^2+\|u\|^2$. Formula~\eqref{eq:phi} intertwines
these real structures. A general pure move need not preserve $\sigma$; that
condition is not discarded. There are also non-unit-fixing norm moves preserving
the compact real form: in standard block order $(A_0,A_1,A_2)$, for
$S\in\mathrm{SL}_3(\C)$,
\[
 (A_0,A_1,A_2)\mapsto
 (SA_0S^\dagger,(S^\dagger)^{-1}A_1,A_2S^{-1})
\]
preserves the norm and the compact reality equations. Determinants and the
cyclic trace prove this directly. Integral unimodular real $S$ preserve the
corresponding entry lattice, but integer matrix entries still do not imply
integer eigenvalues.

\section{A retained ternary quantity on the Leech lattice}

Wilson fixes an octonion basis $1=i_\infty,i_0,\ldots,i_6$ with
$i_0i_1=i_3$ and the cyclic/doubling subscript rules. Let $L$ be the lattice
spanned by the 112 vectors $\pm i_t\pm i_u$ with $t\ne u$, and the 128 vectors
\[
 \tfrac12(\pm1\pm i_0\pm\cdots\pm i_6)
\]
with an odd number of minus signs. Put $s=(-1+i_0+\cdots+i_6)/2$.
At Wilson's normalization, the Leech lattice consists of triples satisfying
\begin{equation}\label{eq:wilson}
 x,y,z\in L,\qquad x+y,x+z,y+z\in L\bar s,
 \qquad x+y+z\in Ls,
\end{equation}
with squared norm $(n(x)+n(y)+n(z))/2$ \cite[\S3, p.~3]{Wilson}.
Both the conjugation bar and the order of right multiplication are retained.
The identification with the Leech lattice is Wilson's theorem, not inferred
here merely from dimension 24.

The linear map
\[
 \iota(x,y,z)=\begin{pmatrix}0&z&\bar y\\\bar z&0&x\\y&\bar x&0\end{pmatrix}
\]
puts this lattice into the three Peirce spaces of the compact Albert algebra.
Its inverse reads the three marked entries. It satisfies
\begin{equation}\label{eq:leechinvariants}
 \Tr(\iota(v)^2)=4\|v\|_\Lambda^2,
 \qquad N_J(\iota(v))=2\Ree((xy)z).
\end{equation}
Thus $\iota/2$ is an isometry for the trace norm. The following basis map
makes the change from Wilson's multiplication convention to the preceding Zorn
convention fully explicit. In compact Zorn coordinates over $\C$, put
\[
 I=(i,0,0,-i),\quad U_j=(0,e_j,-e_j,0),\quad
 V_j=(0,ie_j,ie_j,0).
\]
Then the ordered Wilson basis has images
\begin{equation}\label{eq:wilsonzorn}
 (1,i_0,i_1,i_2,i_3,i_4,i_5,i_6)
 \longmapsto(1,U_1,U_2,I,-U_3,-V_2,-V_3,-V_1).
\end{equation}
For a proof, these eight vectors lie in the compact fixed real form and are
real-linearly independent: their coordinates
$(\Ree\alpha,\operatorname{Im}\alpha,\Ree u_j,\operatorname{Im}u_j)$ form a signed permutation basis.
Each imaginary image squares to $-1$; the ordered products in the seven
cyclic triples $(t,t+1,t+3)$ agree with Wilson's products. The standard-library
certificate checks all 64 basis pairs exactly in $\Q(i)$; bilinearity therefore
proves the multiplication identity for arbitrary octonions. The inverse reads
the displayed real and imaginary coordinates and reverses the signed permutation.
Composing this map entrywise with $\iota$ and $\Phi^{-1}$ puts the same coupled
lattice into the three-block coordinates of Section~5, with its compact real
structure retained. It is not being claimed
that this off-diagonal lattice is closed under the Jordan product or that every
Leech isometry preserves the cubic quantity.

\begin{proposition}[Complete short-shell cubic distribution]\label{prop:leechhist}
On Wilson's 196560 norm-four vectors, the cubic value
$c(v)=2\Ree((xy)z)$ has the following exact distribution:
\[
\begin{array}{c|rrrrrr}
 c&0&\pm1&\pm3&\pm4&\pm5&\pm7\\\hline
 \text{count for each displayed sign}&55248&24288&16800&21504&7392&672.
\end{array}
\]
No other value occurs in this shell. In particular, fixing the quadratic norm
has not fixed the retained ternary Jordan determinant.
\end{proposition}
\begin{proof}
We use Wilson's complete minimal-vector description \cite[\S3, p.~4]{Wilson}.
For each root $\lambda$ of $L$, signed octonion units $j,k$, and all permutations,
the shapes are
\[
 (2\lambda,0,0),\quad
 (\lambda\bar s,\ \pm(\lambda\bar s)j,0),\quad
 ((\lambda s)j,\ \pm\lambda k,\ \pm(\lambda j)k).
\]
The parentheses are part of the formula. The independent checker constructs
all 240 roots, all required products, and the resulting sets; it obtains the
three disjoint cardinalities 720, 11520, and 184320. It tests every vector
against all of \eqref{eq:wilson}, not only the generating shapes, and verifies
its norm and its cubic value with integer arithmetic.

Here are complete executable membership conventions. Represent an octonion by
twice its coordinate vector $v\in\Z^8$. Membership in $L$ is: all coordinates
have one parity $\eps$, and $\sum v_i\equiv2\eps\pmod4$.
The doubled vectors for $s,\bar s$ are respectively
$(-1,1,1,1,1,1,1,1)$ and $(-1,-1,-1,-1,-1,-1,-1,-1)$.
If $B(v,w)$ is the integral bilinear multiplication table, membership of $v/2$
in $L\bar s$ is exactly that $B(v,2s)/4$ is an integral vector satisfying the
$L$ test. Membership in $Ls$ uses $2\bar s$ instead. These formulas follow
from $n(s)=2$ and the octonionic inverse property. For doubled coordinates
$x,y,z$, the cubic value is the scalar coordinate of $B(B(x,y),z)/4$.
These bounded loops and the specified parity tests exhaust precisely the
three displayed sets. Their resulting histogram is the table; the complete
standard-library program, examples of each value, and a sorted-vector hash
are supplied. This is an exhaustive finite certificate with Wilson's shell
identification as its explicit external input.
\end{proof}

\section{Marked Erd\H{o}s--Straus data in a genuine cube}

Fix a positive integer $p$ and an ordered positive integer triple $d=(x,y,z)$. Define
an integral tensor by
\begin{equation}\label{eq:escube}
 \mathcal T_{ijk}=12d_i\delta_{ij}\delta_{ik}-p,
 \qquad i,j,k\in\{1,2,3\}.
\end{equation}
Its trilinear form is
\[
 12\sum_i d_i u_iv_iw_i-p\left(\sum_i u_i\right)
 \left(\sum_i v_i\right)\left(\sum_i w_i\right).
\]
Put $e=(1,1,1)$. Every contraction at the corresponding mark is
\begin{equation}\label{eq:esface}
 M=12\diag(x,y,z)-3p\,ee^T.
\end{equation}

\begin{theorem}[Exact marked cube criterion]\label{thm:esembed}
The encoding \eqref{eq:escube} is injective. Its inverse is
$p=-\mathcal T_{112}$ and $d_i=(\mathcal T_{iii}+p)/12$ on its image.
Moreover,
\[
 \det M=432\big(4xyz-p(xy+xz+yz)\big).
\]
Thus the marked face is singular exactly when $4/p=1/x+1/y+1/z$.
At a passing positive triple it has rank exactly two and kernel
$\Q(1/x,1/y,1/z)$. No global nondegeneracy of the full cube is assumed.
\end{theorem}
\begin{proof}
The inverse reads the off-diagonal constant and then the three diagonal
coefficients, proving injectivity. A determinant expansion of the rank-one
perturbation gives
\[
 \det M=12^3xyz\left(1-\frac p4\left(\frac1x+\frac1y+\frac1z\right)\right),
\]
which is the asserted factor. The diagonal matrix is invertible and a rank-one
change can reduce its rank by at most one. On a passing triple the vector of
reciprocals is in the kernel, so the rank is two and that vector spans it.
\end{proof}

Original selector markings remain attached. For $a=(p+R)/4=hrs$ and
$u=hr^2$, use
\begin{align*}
 E:&\quad R\mid pr+s,\quad \kappa=(pr+s)/R,
 &(a,y,z)&=(hrs,hs\kappa,phr\kappa),\\
 M:&\quad R\mid r+s,\quad \lambda=(r+s)/R,
 &(a,y,z)&=(hrs,phs\lambda,phr\lambda).
\end{align*}
For $0<R<p$ and the retained channel, the divisor inverse is
\[
 u=\frac{a^2}{Ry-pa}\quad(E),\qquad
 u=\frac{pa^2}{Ry-pa}\quad(M).
\]
Let $d_0=\gcd(a,u)$; then $r=u/d_0$, $s=a/d_0$, $h=d_0/r$.
The positive fibre domain $u\mid a^2$ supplies the integrality. The channel,
$p,R$, and raw denominator order are not dropped when using the cube.
The verifier checks these maps on four inherited stress witnesses, including
both channels at 1201 and a positive middle witness at 2521.

\subsection{Why not identify all three arithmetic views?}
All eight invisible additions in Theorem~\ref{thm:glue} can be added to
\eqref{eq:escube} without changing its three marked faces. This is a valid
joint fibre, not a claim of unrestricted dynamics. There is a specific effect
of retaining three \emph{identical} faces: when the resulting cube is
nondegenerate and $M$ has reciprocal kernel $r$, incidence forces
$e\mapsto r\mapsto e\mapsto r$. Hence $H^2(e)=e$. Since $H$ is a translation,
its translation point has order dividing two. This is a consequence of this
particular over-identification, not of all three-slot models.

The following construction avoids it by retaining one arithmetic face and the
actual two other typed views.

\begin{theorem}[An ES mark inside an infinite-return cube]\label{thm:attachment}
For every ordered positive ES solution, its matrix $M$ in \eqref{eq:esface}
is the marked first face of a rational cube projectively equivalent to
\eqref{eq:cube}. After clearing and retaining one scale, the cube is integral.
Its three-slice round trip has infinite order. The complete maps, changes of
basis, and scale are invertible and retain the ES mark.
\end{theorem}
\begin{proof}
Let $A=M_0(1,-1,0)$ for \eqref{eq:cube}. Both $A$ and $M$ have rank two.
Gaussian row and column operations give explicitly invertible matrices
\[
 U_AAV_A=J_2,\qquad U_MMV_M=J_2,\qquad J_2=\diag(1,1,0).
\]
Set $L=U_M^{-1}U_A$ and $R=V_AV_M^{-1}$, so $LAR=M$.
Use
\[
 G_0=\begin{pmatrix}1&0&0\\-2&1&0\\-1&0&1\end{pmatrix},
 \quad G_1=L^T,\quad G_2=R.
\]
Then $G_0e=(1,-1,0)$ and the pullback
\[
 T'(u,v,w)=T(G_0u,G_1v,G_2w)
\]
has first marked face $M$. The inverse applies $G_i^{-1}$ in its own slot.
Multiply all coefficients by their common denominator $b$ and retain $b$;
this changes no projective kernel map. In the new coordinates each kernel
map is $G_j^{-1}\phi_{ij}G_i$, so their full composition is $G_0^{-1}HG_0$.
Theorem~\ref{thm:infinite} proves infinite order. Gaussian pivot choices give
different recorded coordinates, not a false uniqueness claim; with the
deterministic first-nonzero-pivot convention used in the certificate all data
are reproducible.
\end{proof}

For the inherited state
\[
 (p,R,h,r,s,u,\kappa)=(1201,23,17,1,18,17,53),\quad E,
\]
with denominators $(306,16218,1082101)$, the certificate obtains $b=9522$.
The integer cube's first pencil arrays are
\begin{align*}
 &\begin{pmatrix}657018&-35621802&48768510\\
 47783466&-2371784922&-10183692786\\
 -5511721932&284224750146&507029017411\end{pmatrix},\\[1mm]
 &\begin{pmatrix}0&0&1971054\\0&-27363744&1668119022\\
 19044&1824803946&-118273531557\end{pmatrix},\\[1mm]
 &\begin{pmatrix}0&1314036&-85047330\\-82091232&4217974452&8481265998\\
 5477395122&-286083861858&-265144604956\end{pmatrix}.
\end{align*}
Its exact first turn is
\[
 (1,1,1)\to(63653,1201,18)\to(25855837,487577,7866)
 \to(1817,8909,5075).
\]
The first image is the primitive ordered reciprocal triple. The last is the
conjugate of $H(O)$, not a reset of the initial mark. This construction encodes
an existing witness; it is not used circularly as a production theorem for an
unknown prime.

\section{Tetrahedral arithmetic and its exact integral return}

The ES equation also has a direct relation to a genuine tetrahedral cubic,
not only to an $A_4$ action on frames. Let
\[
 \mathscr C:\quad e_3(a_0,a_1,a_2,a_3)=0\subset\mathbb P^3.
\]
This is Cayley's nodal cubic in reciprocal coordinates
\cite[\S1, p.~1]{HB}. Its four coordinate points are singular and are permuted
by $S_4$. Direct substitution gives
\begin{equation}\label{eq:cayleyes}
 e_3(-p,4x,4y,4z)=16\big(4xyz-p(xy+xz+yz)\big).
\end{equation}
The preferred negative coordinate is an arithmetic marking; permutations
forgetting it need not preserve the same prime problem.

On the torus $a_0a_1a_2a_3\ne0$, the reciprocal Cremona map
\[
 [a_i]\longmapsto[b_i],\qquad b_i=\prod_{j\ne i}a_j,
\]
is a bijection between $\mathscr C$ and the plane $\sum b_i=0$ with no zero
coordinate. Its inverse is the same formula. Composing twice gives
$a_i(\prod_ja_j)^2$, the same projective point. Coordinate hyperplanes are the
explicit excluded locus; the four nodes are not treated as ordinary invertible
points of this map.

In the positive ES component the plane coordinates can be oriented as
$(-S,A,B,C)$ with $A,B,C>0$ and $S=A+B+C$. Normalizing by
$\gcd(A,B,C)=1$ gives a unique primitive ordered reciprocal triple.

\begin{proposition}[Complete positive integral return]\label{prop:prime-return}
Let $A,B,C$ be positive coprime as a triple, $S=A+B+C$, and
$L=\lcm(A,B,C)$. Define
\[
 p_0=\frac{4L}{\gcd(4L,S)}.
\]
Then $p_0\ge2$. For a prime $p$, the formulas
\[
 x=\frac{pS}{4A},\qquad y=\frac{pS}{4B},\qquad z=\frac{pS}{4C}
\]
are positive integral ES denominators if and only if $p=p_0$.
All ordered positive ES solutions are obtained exactly once from their
primitive reciprocal triple in this way.
\end{proposition}
\begin{proof}
The reciprocal sum is $4(A+B+C)/(pS)=4/p$. The three integrality conditions
are jointly $4L\mid pS$, equivalently $p_0\mid p$. Since
$S\le3L<4L$, one has $\gcd(4L,S)<4L$ and $p_0\ge2$; primality makes
$p_0\mid p$ equivalent to $p_0=p$. Conversely any ES solution gives the
primitive positive ratio $A:B:C=1/x:1/y:1/z$. Substituting its reciprocal sum
recovers exactly the displayed denominators. Uniqueness follows from primitive
positive normalization.
\end{proof}

For each prime $\ell$, the retained valuation return is the explicit three-slot
piecewise-linear law
\[
 v_\ell(p_0)=\max\{2\,1_{\ell=2}+\max(v_\ell A,v_\ell B,v_\ell C)
                         -v_\ell(A+B+C),\ 0\}.
\]
Thus the multiplicative valuations have a simultaneous max-plus description,
but the valuation of the actual sum remains essential; its arithmetic carries
cannot be replaced by independent slot valuations.

Dividing the coordinates in \eqref{eq:cayleyes} by $4p$ gives
$(-1/4,x/p,y/p,z/p)$. This is an exact occurrence of the marked value $-1/4$
inside ES geometry. No identification with a particular Fable/Jacobian
counterexample point is inferred without a map between those different
constructions.

\section{Exponential coefficients and the retained Terzo period}

The uploaded primary thesis by Giuseppina Terzo \cite{Terzo} uses the free
exponential ring $F=[X,Y]^E$, evaluated at $(X,Y)=(\pi,i)$. Put
\[
 I=\langle Y^2+1,\ E(XY)+1\rangle_E,\qquad A=F/I.
\]
Evaluation $\mathrm{ev}:A\to\C$ is defined unconditionally because the two
relations evaluate to zero. Terzo's Theorem~2.4.1 identifies its kernel as zero
\emph{assuming Schanuel's conjecture}. The separate nonprincipality result is
Theorem~2.4.2 in \S2.4.1. These statements must not be replaced by a casual
``one relation'' count.

All integral polynomial cube, norm, adjugate, and pure-move formulas above base
change to $A$. Evaluation commutes with them coefficient by coefficient. If
$J=\ker(\mathrm{ev})$, then the kernel on the free cube module is exactly
$J^{27}$. This is an unconditional equality, not an assertion that $J$ vanishes.
Under Schanuel, Terzo's theorem makes the evaluation faithful. The Jordan
product further requires localization at two as an underlying ring; no
canonical extension of exponentiation to this localization is asserted.

This coefficient layer retains exponential terms in $\pi,i$, including nested
terms, rather than pretending that they span a three-dimensional vector space
over $\Q$. Three slots and coefficient transcendence are independent parts of
the specification. For example, powers of any transcendental number are
$\Q$-linearly independent, so a finite-dimensional $\Q$-algebra closed under
multiplication cannot contain that number. This elementary fact does not
prevent a rank-three or three-slot construction over the exponential
coefficient ring; it specifies which notion of linearity is being used.

There is also an actual all-integer link to Section~\ref{sec:moore}. Put
$\varpi=2\bar X\bar Y\in A$. Then
\[
 E(\varpi)=E(\bar X\bar Y)^2=1,\qquad
 \mathrm{ev}(\varpi)=2\pi i.
\]
The subgroup $\Z\varpi$ is infinite cyclic unconditionally: evaluating
$n\varpi=0$ forces $n=0$. Therefore
\begin{equation}\label{eq:periodmap}
 \Z\varpi\longrightarrow\langle H\rangle,\qquad
 n\varpi\longmapsto H^n
\end{equation}
is an explicit isomorphism of marked groups by Theorem~\ref{thm:infinite}.
It also identifies the actual period lattice $2\pi i\Z$ with these retained
cube returns. Under Schanuel only, the complete exponential kernel in $A$
equals $\Z\varpi$: evaluate a kernel element into $2\pi i\Z$ and use
injectivity. The unconditional marked subgroup map \eqref{eq:periodmap} does
not depend on that strengthening. It is a group/lattice map, not a map of
exponential rings, an analytic zero-localization theorem, or an ES proof.

Quintic algebraic coefficients can similarly be retained by base change to
$k[\theta]/(f(\theta))$. This does not identify tensor arity with polynomial
degree, or prove a specific link to a proposed quintic core. Neither that core
nor a Schanuel-strength transcendental classification is needed for the
unconditional cube and graded-operator theorems.

\section{Executable scope and the resulting pointwise target}

The package provides three standard-library programs:
\begin{verbatim}
python verify.py --out certificate.json
python -O verify.py
python graded_lie.py
python -O graded_lie.py
python leech_ternary.py
python -O leech_ternary.py
\end{verbatim}
Normal and optimized outputs are checked for byte equality separately for each
program. No removable Python assertion is used. The first program includes
formal sparse-polynomial identities, exact rational kernel calculations,
complete reduction point sets, a 27-cell change-of-coordinate inverse, and
marked ES stress witnesses. The second verifies all 3003 basis commutators in
a 78-dimensional three-graded matrix algebra. The third reconstructs and tests
all 196560 listed Wilson minimal vectors. Imported mathematical theorems,
formal polynomial checks, and finite exhaustive checks retain different scopes.
No Lean run or independent reviewer is claimed; a declaration-level Lean plan
is supplied.

The principal construction is no longer confined to a frame rotation. It has
three typed views of one tensor, nontrivial all-integer return, an action mixing
all three exceptional blocks, an explicit coupled Leech sublattice, and a
reversible ES arithmetic mark. The norm-preserving moves genuinely change the
quadratic trace, so the earlier frame-only invariant is not an obstruction to
this enlarged class.

What has not been proved is that these transformations produce, for every
required prime, a positive primitive triple whose exact return value in
Proposition~\ref{prop:prime-return} is that prime. In particular, an intrinsic
split-Albert equation, a nonempty complex orbit, or a rational point in a
continuous coordinate chart is not silently converted to a point in that
positive integral fibre. A precise next theorem is a terminating three-slot
integral operation with $p_0=p$, or an equivalent marked shell-production
operation, while preserving all divisor exponents and signs. The present
construction provides the explicit maps and invariant that such an operation
must change. It neither proves global ES nor gives an obstruction to the
whole proposed programme.

\begin{thebibliography}{9}
\bibitem{BH}
Bhargava, M., \& Ho, W. (2013). \emph{Coregular spaces and genus one curves}.
 arXiv:1306.4424v1. Used: \S\S3.2.2--3.2.3, printed pp.~19--20;
 Theorem~5.1 and \S5.1.1, pp.~32--35, especially the paragraph following
 Lemma~5.3. \url{https://arxiv.org/abs/1306.4424}.
\bibitem{BP}
Baez, J. C., \& Pumpl\"un, S. (2026). \emph{Three-dimensional geometry in
 exceptional mathematics}. arXiv:2609.07538v1, pinned HTML title (the moving
 abstract title may differ). Used: \S5, equations (10)--(15), Theorem~11 and
 Appendix A; \S6, equations (16)--(20). The first-Tits and Br\"uhne antecedents
 are explicitly credited there. \url{https://arxiv.org/html/2609.07538v1}.
\bibitem{Milne}
Milne, J. S. (2006). \emph{Elliptic curves}. BookSurge.
 Used: Chapter II, Theorem~4.1 and Corollary~4.2, pp.~62--64, and
 Corollary~5.7, p.~66. Author-hosted edition:
 \url{https://www.jmilne.org/math/Books/ectext6.pdf}.
\bibitem{Wilson}
Wilson, R. A. (2009). Octonions and the Leech lattice.
 \emph{Journal of Algebra, 322}(6), 2186--2190.
 doi:10.1016/j.jalgebra.2009.03.021. Used: revised author preprint,
 20 March 2009, \S\S2--4, printed pp.~2--5; the conjugation bars and
 minimal-vector parentheses were inspected on page images.
 \url{https://maths.qmul.ac.uk/~raw/pubs_files/octoLeech1rev.pdf}.
\bibitem{HB}
Heath-Brown, D. R. (2002). \emph{The density of rational points on Cayley's
 cubic surface}. arXiv:math/0210333v1. Used only for the classical Cayley
 cubic identification, \S1, p.~1; no counting theorem is invoked.
 \url{https://arxiv.org/abs/math/0210333}.
\bibitem{Terzo}
Terzo, G. (n.d.). \emph{Consequences of Schanuel's conjecture in exponential
 algebra} [Thesis; user-supplied primary document]. Used: \S2.4,
 Theorem~2.4.1, and \S2.4.1, Theorem~2.4.2; coefficient and period scope
 distinguished in the text. No publication year is inferred from the upload.
\end{thebibliography}
\end{document}
